\cotefile{8}
\scanpage{35}
% archive sheet number: 34
\mspage{84}
\hfill n\textsuperscript{o} 382

application de classe $C^{r-k}$ de $f^{-1}P^k(X',\pi')=X\times_{X'} P^k(X',\pi')$ dans $P^k(X,\pi)$, qui, si $X$, $X'$, $Y'$ sont des variétés pures, est un morphisme pour les fibrations sur $X$.

7.2.7. Soient $X$, $Y$ des variétés de classe $C^r$, $\varphi:X\longrightarrow Y$ une submersion, $E$ un fibré vectoriel de classe $C^r$ sur $Y$, et $E'$ le fibré vectoriel image réciproque de $E$ suivant $\varphi$. Soit $k$ un entier fini $\leq r$ et soit $P^k(Y,E)_{\unsure{o|0}}$ le fibré vectoriel sur $Y$ dont la fibre $P^k_y(Y,E)_o$ en $y\in Y$ est l'espace des jets de sections nulles en $y$. Soient $x\in X$ et $y=f(x)$. On a un morphisme canonique de $P^k_x(X)\otimes\varphi^{-1}P^k_y(Y,E)_o$ dans $P^k_x(X,E')$, dont l'image est l'ensemble des jets de section de $E'$ en $x$ dont la restriction à $X_y=\varphi^{-1}(y)$ est nulle, et une suite exacte de morphismes d'espaces fibrés vectoriels
\[
P^k(X)\otimes\varphi^{-1}(P^k(Y,E)_o)\longrightarrow P^k(X,E')\longrightarrow P^k(X/Y,E)\longrightarrow 0\;.
\]
La fibre de $P^k(X/Y,E)$ située sur $x$ s'identifie à $J^k_x(X_y,E_y)$, où $E_y$ est la fibre de $E$ sur $y$.

Soient encore $Y'$ une variété pure de classe $C^k$, $f:Y'\longrightarrow Y$ un morphisme de $Y'$ dans $Y$, et $X'=Y'\times_Y X$ le produit de $X$ et $Y'$ au-dessus de $Y$. Alors l'application canonique $\pi':X'\longrightarrow Y'$ est une submersion ($\hspace{3.6em}$) et l'on a
\[
P^k(X'/Y',f^{-1}E)=P^k(X/Y,E)\otimes_{Y'}X'\;.
\]

7.2.8. Soient $X$, $Y$ des variétés de classe $C^r$, $\pi:Y\longrightarrow X$ une submersion, $r'\leq r$ et $k$ un entier fini $\leq r'$, et soient $Y'$, $X'$ les variétés de classe $C^{r'}$ obtenues à partir de $Y$, $X$ par affaiblissement de structure ($\hspace{3.5em}$). Alors $P^k(X',\pi)$ est obtenue par affaiblissement de structure de $C^{r-k}$ à $C^{r'-k}$ à partir de $P^k(X,\pi)$.
\scanpage{36}
% archive sheet number: 35
On a $A\otimes_{\mathbb Z_p}W$ \unsure{étale} sur $A$, ext. \unsure{résiduelle|résiduels}
$\begin{aligned}[t] k\otimes_{\mathbb F_p}k&\simeq k^\Sigma,\quad\Sigma=\operatorname{Aut} k\\[-2pt]&\hphantom{\simeq k^\Sigma,\quad\Sigma}=\operatorname{Aut} W\end{aligned}$
\[
A\otimes_{\mathbb Z_p}W\simeq A^\Sigma,\qquad \Sigma\simeq\mathbb Z/f\mathbb Z\quad\text{en correspondance}
\]
\hfill bijection canonique avec

\hfill $[0,f-1]$ par $\unsure{\nu|v}\longmapsto\sigma^\nu$

\hfill \deleted{$(\unsure{\sigma^\nu=\mathrm F^p|\sigma=\mathrm F^p})$}

On a
\[
\sigma_A(x_0,x_1,\ldots,x_{f-1})=(x_1,\ldots,x_{f-1},x_0)
\]

La donnée d'un module \inserted{libre} $M$ \inserted{sur $A\otimes_{\mathbb Z_p}W$} équivaut à la donnée de modules $(\deleted{\unsure{A}}M_\nu)_{\nu\in\Sigma}$ sur $A$ ; le rang de $\deleted{\unsure{A}}M$ est la somme des rangs des $M_\nu$. \inserted{(on le veut égal à $h$)} La donnée d'un \unsure{endom} $F_M$ $\sigma_A$-linéaire équivaut à la donnée d'hom. $A$-linéaires
\diagram{\begin{tikzcd}[ampersand replacement=\&,column sep=small]
M_0\arrow[r,"\varphi_0"] \& M_1\arrow[r,"\varphi_1"] \& \cdots\arrow[r] \& M_{f-1}\arrow[r,dash,"\varphi_{f-1}"] \& M_0
\end{tikzcd}}

La condition $(*)$ \deleted{\unsure{s'écrit}} implique que les $f$ composés circulaires sont des isogénies, donc les $\varphi_i$ tous injectifs, donc les $\varphi_i$ des isogénies (car les rangs des $M_\nu$ tous égaux \ldots). \deleted{Donc} --- Utilisant maintenant le fait que la somme des rangs des $M_\nu$ est $h$, on trouve que les $M_\nu$ sont inversibles.
% slanted
\margin{ce serait aussi une \deleted{conséquence de $(*1)$} \deleted{\unsure{cas}} de $(*)$ avec $\pi$ un élément non nul \uline{quelconque} de $A$}

\deleted{On peut alors trouver un unique module inversible}

On peut alors identifier tous les $M_\nu$ à \deleted{$M_0=\Omega$,} \deleted{$\Omega$ sur $A$,} $M_0=t$, \deleted{\unsure{qu'on a|qui a i}} en choisissant \deleted{\unsure{comme|can.}} pour \unsure{tt|$M$} $\unsure{\nu|\nu'}$,
\inserted{$0\leq\nu\leq f-1$,} l'isom $M_0\xrightarrow{\sim}M\unsure{.|\!{}_\nu|\!{}_{i}}$ donné par
$\deleted{\unsure{p}}\pi^{-m_\nu}(\varphi_{\nu-1}\varphi_{\nu-2}\ldots\varphi_0)$, où $m_\nu$ est un entier $\geq 0$ uniquement déterminé ($\unsure{0\leq|0=}m_0\leq m_1\leq\ldots\leq m_{f-1}$). On trouve donc un isom can.
\[
 M\simeq\Omega\otimes_{\mathbb Z_p} W=\Omega\otimes_A(A\otimes_{\mathbb Z_p}W)
\]
\scanpage{37}
% archive sheet number: 36
\mspage{83}
\hfill n\textsuperscript{o} 382

des germes en $x$ de sections qui ont un contact d'ordre $\geq k'$ avec la section nulle ; cet espace est la fibre en $x$ d'un sous-fibré vectoriel de $P^k(X,E)$, qui est aussi un module sur $P^k(X)$. En particulier, le noyau de $P^k(X,E)\longrightarrow P^{k-1}(X,E)$ s'identifie au produit tensoriel $S^k(T(X)^*)\otimes E$ de $E$ avec la $k$-ième puissance symétrique du dual du fibré tangent à $X$.

7.2.4. Soient $X$, $Y$ des variétés de classe $C^r$, $\pi:Y\longrightarrow X$ une submersion, $k'\leq k$ des entiers finis $\leq r$, $\handwritten{\ell}=k-k'$, et soit $\pi_1$ la projection de $P^{k'}(X,\pi)$ sur $X$. L'application canonique de $J^k(X,Y)$ dans $J^{\handwritten{\ell}}(X,J^{k'}(X,Y))$ définie au n\textsuperscript{o} 7.1.10 induit une application de classe $C^{r-k}$ de $P^k(X,\pi)$ dans $P^{\handwritten{\ell}}(X,\pi_1)$ qui, lorsque $K$ est de caractéristique zéro, est un isomorphisme de $P^k(X,\pi)$ sur une sous-variété de $P^{\handwritten{\ell}}(X,\pi_1)$ ; cette application est un morphisme de fibrés sur $X$ si $X$ et $Y$ sont des variétés pures.

7.2.5. Soient $X$, $Y$, $Y'$ des variétés de classe $C^r$, $\pi:Y\longrightarrow X$ et $\pi':Y'\longrightarrow X$ des submersions, $g:Y\longrightarrow Y'$ un morphisme de $Y$ dans $Y'$ tel que $\pi'\circ g=\pi$, \handwritten{\textsuperscript{(1)}}%
% The matching (1) and its text stand at the foot of the sheet.
\inserted{\handwritten{(1) et $k$ un entier $\leq r$. L'application $j^k(\mathrm{Id},g)$}} (cf. n\textsuperscript{o} 7.1.11) induit une application $P^k(g):P^k(X,\pi)\longrightarrow P^k(X,\pi')$ de classe $C^{r-k}$ qui, lorsque $X$, $Y$, $Y'$ sont des variétés pures, est un morphisme pour les fibrations sur $X$. Si $Y''$ est une variété de classe $C^r$, $\pi'':Y''\longrightarrow X$ une submersion et $g':Y'\longrightarrow Y''$ un morphisme tel que $\pi'=\pi''\circ g'$, alors $P^k(g'\circ g)=P^k(g')\circ P^k(g)$. Si $E\xrightarrow{\lambda} F\xrightarrow{\mu} G$ est une suite exacte de fibrés vectoriels sur $X$, alors
\diagram{\begin{tikzcd}[ampersand replacement=\&]
P^k(X,E)\arrow[r,"P^k(\lambda)"] \& P^k(X,F)\arrow[r,"P^k(\mu)"] \& P^k(X,G)
\end{tikzcd}}
est une suite exacte. Si $A$ et $B$ sont des fibrés vectoriels sur $X$, alors
\[
P^k_x(X,A\otimes B)\cong P^k_x(X,A)\otimes_{P^k_x(X)} P^k_x(X,B)\;,\quad\text{pour tout }x\in X\;.
\]

7.2.6. Soient $X$, $X'$, $Y'$ des variétés de classe $C^r$, $f:X\longrightarrow X'$ un morphisme, $\pi:Y'\longrightarrow X'$ une submersion, $Y=X\times_{X'}Y'$ (n\textsuperscript{o} 4.11.3) et $\pi$ la submersion canonique de $Y$ sur $X$, ($\hspace{1.5em}$). L'application $j^k(f,\mathrm{Id}.)$ (cf. n\textsuperscript{o} 7.1.11) induit une
\scanpage{38}
% archive sheet number: 37
\deleted{\unsure{Et} \inserted{Alors} $F_M$ est donné par les entiers $m_i$ \inserted{et par $\varphi_{f-1}$} \unsure{mais qui} vérifient donc \unsure{$=$}, les $\unsure{m_i,m_{i-2}|m_i,m_{i-1}|m_{i+1}-m_i}$}

On a choisi les isom de telle façon que chaque $\varphi_{\unsure{\nu|v}}$ ($0\leq\nu\leq f-2$) soit de la forme
\[
\varphi_\nu=\pi^{r_{\deleted{\unsure{0}}\nu}}\mathrm{Id}_\Omega
\qquad (r_\nu\geq 0\deleted{\unsure{)}},\;0\leq\nu\leq f-2)
\]
et $\varphi_{f-1}$ est donné alors par
\[
\varphi_{f-1}=\lambda\,\mathrm{id}_\Omega \qquad \lambda\in A,\;\inserted{\text{s'écrit}}\ \lambda=\pi^{r_{f-1}}u,\;u\in A^*
\]

\deleted{\unsure{La co}} On a alors
\[
F_M^f=\pi^{r_0+\cdots+r_{f-2}+r_{f-1}}u=\varpi\qquad u\in A^*
\]

\deleted{On aura}
\[
\deleted{M/F_M M\simeq A/\pi^{r_{f-1}}A+A/\pi^{r_{\unsure{0|i}}}A+\ldots+A/\pi^{r_{f-2}}A}
\]

La condition $(*)'$ s'écrit

\deleted{\unsure{donc} \unsure{donc|$d_K$} $\pi$}

\[
\begin{aligned}
&p\ \text{divisible par }\pi^{r_\nu}\ \text{pour tt }0\leq\nu\leq f-2,\\
&\hphantom{p\ \text{divisible }}\text{par }u\pi^{r_{f-1}}
\end{aligned}
\]
i.e.
\[
\left\{\begin{aligned}
&\boxed{\deleted{\unsure{0\leq}}r_\nu\leq e}\quad\text{pour }0\leq\nu\leq f-1\\
&\qquad\boxed{\deleted{\unsure{r_{f-1}+v(u)\leq e}}}
\end{aligned}\right.
\]
On a
\[
M/F_M M=A/u\pi^{r_{f-1}} A\unsure{+|-}A/\pi^{r_0}A+\ldots+A/\pi^{r_{f-2}}A
\]
est de \deleted{\unsure{dimension}} longueur
\[
\begin{aligned}
&\boxed{\dim G_0=\deleted{\unsure{v(u)+}} \inserted{\unsure{\operatorname{long} A(\varpi A)\mathrel{\overset{?}{=}}|\operatorname{long} A/(\varpi A)\mathrel{\overset{?}{=}}}}\ r_0+\ldots+r_{f-2}+\unsure{(}r_{f-1}\deleted{\unsure{+v(u)}})}\\
&\hspace{7em}(\leq fe=h)
\end{aligned}
\]

\deleted{Supposons maintenant}

Donc on trouve
\[
\boxed{\begin{aligned}
\dim G_0=1&\Longleftrightarrow F_M^f\ \deleted{\unsure{\text{une}}}\inserted{=\varpi_M,\;\varpi\ \text{une}}\ \text{\uline{uniformisante} de }A\\[-2pt]
&\qquad\text{(qu'on peut supposer alors }=\pi,\ \text{quitte à changer }\pi\,!\,)\\
&\Longleftrightarrow\text{tous les }r_\nu\text{ sauf un seul }\inserted{r_{\nu_0}}\text{ sont nuls},\\
&\qquad r_{\nu_0}=1,\quad u\text{ une unité}\\
&\hspace{9em}(u=1\text{ si on veut }\varpi=\pi)
\end{aligned}}
\]

\unsure{Ce|G}
\scanpage{39}
% archive sheet number: 38
\mspage{82}
\hfill n\textsuperscript{o} 382

variétés sont pures, et si $X'=X$, (resp. $Y'=Y$), l'application $j^k(\mathrm{Id}.,g)$ (resp. $j^k(f,\mathrm{Id}.)$) est un morphisme pour les fibrations sur $X$ (resp. sur $Y$).

7.1.12. Soient $X$, $Y$ des variétés de classe $C^r$, $r'\leq r$ et $k$ un entier fini $\leq r'$. Soient $X'$ et $Y'$ les variétés de classe $C^{r'}$ obtenues à partir de $X$ et $Y$ par affaiblissement de structure (n\textsuperscript{o} 4.13.1). Alors $J^k(X',Y')$ est obtenue par affaiblissement de structure de $C^{r-k}$ à $C^{r'-k}$ à partir de $J^k(X,Y)$.

7.2 - \uline{Jets de sections}.

7.2.1. Soient $X$, $Y$ des variétés de classe $C^r$, $\pi:Y\longrightarrow X$ une submersion et $k\leq r$. On note $P^k_x(X,\pi)$ l'ensemble des éléments $f$ de $J^k_x(X,Y)$ tels que $j^k_y(\pi)\circ f$ soit le jet de l'application identique de $X$, $y$ étant le but de $f$. Le jet $f$ a cette propriété si et seulement si c'est le jet d'ordre $k$ d'un germe en $x$ d'application $f$ de $X$ dans $Y$ tel que $\pi\circ f$ soit le germe de l'application identique.

7.2.2. On suppose dorénavant $k$ fini. L'ensemble $P^k_x(X,\pi)$ est une sous-variété fermée de $J^k_x(X,Y)$, et la réunion $P^k(X,\pi)$ des $P^k_x(X,\pi)$ ($x\in X$) est une sous-variété fermée de $J^k(X,Y)$. Si $X$ et $Y$ sont des variétés pures, les fibrations de $J^k(X,Y)$ sur $X\times Y$, $X$ et $Y$ (cf. n\textsuperscript{o} 7.1.6) induisent des fibrations de $P^k(X,\pi)$, de bases respectives $\pi(X)\times Y$, $\pi(X)$ et $Y$. Si $\pi$ est la projection d'une fibration $\xi$, on écrira aussi $P^k(X,\xi)$ pour $P^k(X,\pi)$. Si $\pi$ est la projection d'un fibré vectoriel $E$, alors $P^k(X,\pi)$ est un fibré vectoriel sur $X$, et un fibré en modules sur le fibré en algèbres $P^k(X)$ ; on le notera aussi $P^k(X,E)$. Si de plus $E$ est le fibré vectoriel trivial $F_X$ de fibre $F$ sur $X$, alors
\[
P^k(X,\pi)\cong P^k(X)\otimes F_X\cong J^k(X,F)\;.
\]

7.2.3. Soit encore $k'$ un entier $\leq k$. L'application $J^k(X,Y)\longrightarrow J^{k'}(X,Y)$ (cf. n\textsuperscript{o} 7.1.2) induit une application $P^k(X,\pi)\longrightarrow P^{k'}(X,\pi)$ qui est un morphisme pour les fibrations sur $X$. Si $Y=E$ est un fibré vectoriel sur $X$, le noyau du morphisme $P^k_x(X,E)\longrightarrow P^{k'}_x(X,E)$ ($x\in X$) s'identifie à l'espace vectoriel des jets
\scanpage{40}
% archive sheet number: 39
Donc \deleted{\unsure{la|les} \illegible{1}} on trouve

\uline{Proposition} $A$, \inserted{$h$,} $W,e,f$, \inserted{$h=ef$,} comme plus haut, $G_0$ un groupe de B.T: \inserted{de hauteur $h$} sur \deleted{$A$} $k$ sur lequel $A$ opère, \deleted{Alors \unsure{$F$|si}} $M$ son module de Dieudonné (module libre de rang $h$ sur $W$). Alors $F_{G_0}^f$ est de la forme $\varpi_{G_0}$, où $\varpi\in A$, et on a $\dim G_0=\operatorname{long} A/\varpi A=v(\varpi)$.

Alors \deleted{conditions équivalentes}

Conditions équivalentes

(i) $\dim G_0=1$

(ii) $F_{G_0}^f=\pi_{G_0}$, $\pi$ une \uline{uniformisante} de $A$ (uniqt déterminée !)

(iii) Le module de Dieudonné $M$ de $G_0$ \inserted{(sur $W$)}, avec $F_M$ et action de $A$, est donné par un triple $(\Omega,\unsure{\nu_0|\nu},\pi)$, où $\Omega$ est un module inv. sur $A$, $0\leq\nu\leq f-1$, $\pi$ une uniformisante de $A$ par
\[
\left\{\begin{aligned}
M&=\Omega\otimes_{\mathbb Z_p}W=\Omega\otimes_A(A\otimes_{\mathbb Z_p}W)\qquad(\text{isom de }A\otimes_{\mathbb Z_p}W\text{-mod})\\
F_M&=\mathrm{id}_\Omega\otimes_{\mathbb Z_p}\sigma+\operatorname{pr}_{\nu}[(\pi-1)\otimes\sigma]\deleted{\unsure{\operatorname{pr}_{\nu}}}
\end{aligned}\right.
\]
où $\operatorname{pr}_\nu$ est la \deleted{facteur} projection de $M$ sur son $\nu$-ième facteur \inserted{$M_\nu$} (correspondant à la décomposition $A\otimes_{\mathbb Z_p}W\simeq A^{[0,f-1]}$)

% slanted
\margin{\uline{NB} Tous ces groupes (pour $\nu$ \unsure{variable}) sont \unsure{déduits|isomorphes}, de l'un d'entre eux (disons pour $\nu=0$) par \unsure{transformation|transport|transformations} à l'aide des $f$ automorphismes $\sigma^i$ de $k=\mathbb F_q=\mathbb F_{p^f}$.}

\uline{NB} Revenons au cas général où on ne suppose pas néc. $\dim G_0=1$ \inserted{(mais $G_0$ de hauteur $h$ $\ldots$)}. Soit $G'_0$ défini comme $G_0$ par $\Omega'$, \deleted{\illegible{1}} les $r'_\nu$ $(0\leq\nu\leq f-1)$ $u'\in A^*$. On cherche les hom compatibles avec l'action de $A$ de $G_0$ dans $G'_0$, i.e. les hom de $A\otimes_{\mathbb Z_p}W$-modules de $M'$ en $M$, compatibles avec $F_M,F_{M'}$. On trouve \unsure{facilt|facilemt|facilement} que pour qu'il existe un tel hom qui \deleted{\unsure{soit}} $\ne 0$ ($\Longleftrightarrow$ une isogénie !), il f. et suffit qu'on ait $\varpi=\varpi'$ i.e. $F_{G_0}^f=F_{G'_0}^f$ (dans $A$). Dans ce cas, \deleted{\unsure{les}} le gpe des hom en question, comme $A$-module, est isom au \deleted{\unsure{$E'$}} sous-module $H\subset\operatorname{Hom}(\Omega',\Omega)$ formé des $v_0:\Omega'\to\Omega$ tels qu'on ait
\[
v_0\pi^{(\sum_0^\nu r_i)-(\sum_0^\nu r'_i)}\in\operatorname{Hom}(\Omega',\Omega)\quad\inserted{\text{pour tt }0\leq\nu\leq f-2},\quad\text{i.e.}\quad \operatorname{val}(v_0)+\sum_0^\nu r_i-\sum_0^\nu r'_i\geq 0
\]
\scanpage{41}
% archive sheet number: 40
\mspage{81}
\hfill n\textsuperscript{o}382

$f\in J^k(X,K^q)_{\unsure{o|0}})$.

7.1.7. Si $Y=K$, la structure d'algèbre des $P^k_x(X)$ munit $P^k(X)$ d'une structure d'espace fibré vectoriel en algèbres de base $X$, de classe $C^{r-k}$.

7.1.8. Soit $X$ une variété de classe $C^r$, et soit $x\in X$. L'espace vectoriel $J^{1}_{o}(K,X)_x$ s'identifie canoniquement à l'espace tangent $T_x(X)$ à $X$ en $x$. La réunion $T(X)=\handwritten{\bigcup}_{x\in X}T_x(X)$ des espaces tangents aux points de $X$ s'identifie canoniquement à $J^{1}(K,X)$, et admet donc une structure d'espace fibré vectoriel sur $X$, de classe $C^{r-1}$. C'est le \uline{fibré tangent} $T(X)$ à $X$.

7.1.9. Soient $X$ et $Y$ des variétés de classe $C^r$ et $k$ un entier fini $\leq r$. Un morphisme $f:X\longrightarrow Y$ définit une application $x\longrightarrow j^k_x(f)$ de $X$ dans $J^k(X,Y)$ qui est une section de classe $C^{r-k}$ pour la fibration de $J^k(X,Y)$ sur $X$. L'application $J^k(X,Y)\longrightarrow J^{k'}(X,Y)$ ($k'$ entier $\leq k$, cf. n\textsuperscript{o} 7.1.2) est un morphisme pour les fibrations sur $X$ et pour les fibrations sur $Y$.

7.1.10. Soient $X$, $Y$ des variétés de classe $C^r$, $k'\leq k$ des entiers finis $\leq r$, $\handwritten{\ell}=k-k'$. Soit $f$ une application de classe $C^r$, dans $Y$, d'un ouvert $U$ de $X$. l'application $x\longrightarrow j^{k'}(f)$ de $U$ dans $J^{k'}(X,Y)$ est de classe $C^{r-k'}$ et son jet d'ordre $\handwritten{\ell}$ en $x$ ne dépend que de $j^k(f)$. On obtient ainsi une application canonique de classe $C^{r-k}$ de $J^k(X,Y)$ dans $J^{\handwritten{\ell}}(X,J^{k'}(X,Y))$ qui, si $K$ est de caractéristique zéro, est un isomorphisme de $J^k(X,Y)$ sur une sous-variété de $J^{\handwritten{\ell}}(X,J^{k'}(X,Y))$ et qui, si $X$ et $Y$ sont des variétés pures, est un morphisme pour les fibrations sur $X$.

7.1.11. Soient $X$, $Y$, $X'$, $Y'$ des variétés de classe $C^r$, $f:X\longrightarrow X'$ et $g:Y'\longrightarrow Y$ des morphismes, $x\in X$, $y'\in Y'$, $x'=f(x)$, $y=g(y')$. Les morphismes $f$ et $g$ définissent une application $j^k_{x'}(f,g)_y:J^k_{x'}(X',Y')_{y'}\longrightarrow J^k_x(X,Y)_y$ caractérisée par $u\longrightarrow P^k_{y'}\circ u\circ P^k_x(f)$. On en déduit une application, $j^k(f,g)$ de $f^{-1}J^k(X',Y')=X\times_{X'}J^k(X',Y')$ dans $J^k(X,Y)$, de classe $C^{r-k}$. Si ces
\scanpage{42}
% archive sheet number: 41
En particulier, si $G_0,G'_0$ satisfont aux conditions de la proposition, alors $\operatorname{Hom}(G_0,G'_0)\ne 0$, i.e. $G_0$ et $G'_0$ sont isogènes, ssi $\pi=\pi'$. \deleted{\unsure{Soit}} S'il en est ainsi, alors \deleted{\illegible{2}} on a
\[
\left\{\begin{aligned}
&\text{si }\nu\unsure{\leq|<}\nu'\quad\deleted{G_0=G'_0},\quad \operatorname{Hom}_A(\unsure{G_0|\Omega_0},G'_0)\xrightarrow{\sim}\operatorname{Hom}(\Omega',\Omega)\\
&\text{si }\nu>\nu'\qquad\operatorname{Hom}_A(G_0,G'_0)\simeq\pi\operatorname{Hom}(\Omega',\Omega)
\end{aligned}\right.
\]

\uline{Re NB} Revenons au cas général où on ne suppose pas néc. $\dim G_0=1$, \uuline{ni} $\varpi$ hauteur $=h$. Classifions à isogénie près \deleted{\unsure{les}} ces $G_0$, donc les $M_\nu$ $(0\leq\nu\leq f-1)$ deviennent des vectoriels \deleted{de \unsure{dim.}} \inserted{sur le} corps \deleted{\unsure{$A$}}$K$, isomorphes par les $\varphi_\nu$ $(0\leq\nu\leq f-2)$. Alors \deleted{\unsure{il est clair encore que}}
\[
M\simeq\deleted{\unsure{M}}\Omega\otimes_{\deleted{\unsure{\mathbb Z_p}}\mathbb Q_p}W\simeq\Omega\otimes_{\unsure{K_p|K}}\bigl(\deleted{\unsure{k}}K\deleted{\unsure{\otimes}}\otimes_{\mathbb Q_p}\underset{\inserted{\text{corps des fractions de } W}}{K_0}\bigr).
\]

et $F$ est déterminé par la connaissance de $\varphi_{f-1}$, qui est un endom de $M$, soit $\varpi$, de sorte qu'on a
\[
F_M^{\unsure{f|s}}=\varpi\otimes_{\mathbb Q_p}\mathrm{id}_{\deleted{\unsure{k}}K_0}=\varpi\otimes_K(K\otimes_{\mathbb Q_p}K_0)
\]
La condition \textquotedblleft d'effectivité\textquotedblright\ en termes de $\varpi$ est
\[
\left\{\begin{aligned}
&\varpi\ \text{laisse invariant un réseau}\\
&p^f\varpi^{-1}\ \text{laisse invariant un réseau}
\end{aligned}\right.
\]
\scanpage{43}
% archive sheet number: 42
\mspage{79}
\hfill n\textsuperscript{o} 382

$|\alpha|\leq k$ (cf. n\textsuperscript{o} 2.5.3 et 3.5.2) ; si $r\ne\omega$, ces valeurs peuvent être choisies arbitrairement.

L'ensemble $J^k(K^p,K^q)$ s'identifie canoniquement à $K^p\times K^q\times L^k_{p,q}$, où l'on a posé $L^k_{p,q}=J^k_{\unsure{o|0}}(K^p,K^q)_o$. L'ensemble $L^k_{p,q}$ admet une structure d'espace vectoriel sur $K$, caractérisée par
\[
j^k_o(f)+j^k_o(g)=j^k_o(f+g)\;,\quad\alpha.j^k_o(f)=j^k_o(\alpha.f)\qquad(\alpha\in K)\;,
\]
$f$ et $g$ étant des applications dans $K^q$ d'un voisinage ouvert de l'origine de $K^p$, de classe $C^r$, envoyant l'origine sur l'origine.

Supposons $k$ fini. Alors $L^k_{p,q}$ s'identifie canoniquement à $K^{q,c(k,p)}$, avec $c(k,p)=\binom{k+p}{p}$, sa base canonique étant formée des parties principales d'ordre $k$ des applications polynomiales $p_{j,\handwritten{\beta}}:x\longrightarrow(f^1(x),\ldots,f^q(x))$, où $f^i(x)=\handwritten{\delta}^i_j.x^{\handwritten{\beta}}$ ($j=1,\ldots,q$ ; $|\handwritten{\beta}|\leq k$). Cette identification est compatible avec la structure vectorielle définie précédemment. Les espaces $L^k_{p,q}$ et $J^k(K^p,K^q)$ sont ainsi munis canoniquement d'une structure de variété analytique.

L'ensemble $L^k_p$ des jets inversibles d'ordre $k$ de $K^p$ dans $K^p$, de source et but l'origine, muni de la loi de composition des jets, est une variété de groupe de classe $C^\omega$ ; le groupe $L^k_p\times L^k_q$ opère analytiquement à gauche sur $L^k_{p,q}$ par la composition des jets :
\[
(a,b).c=a\circ c\circ b^{-1}\qquad (a\in L^k_p,\ b\in L^k_q,\ c\in L^k_{p,q})\;.
\]

Soit $k'\leq k$. L'application naturelle $L^k_{p,q}\longrightarrow L^{k'}_{p,q}$ (cf. 7.1.4) est surjective, linéaire. Son noyau est l'ensemble des jets d'ordre $k$ de $K^p$ dans $K^q$ qui ont un contact d'ordre $\geq k'$ avec l'application nulle. L'application naturelle $L^k_p\longrightarrow L^{k'}_p$ est un morphisme surjectif de variétés de groupe de classe $C^\omega$.

7.1.6. Soient $X$ et $Y$ des variétés de classe $C^r$, $k$ un entier fini $\leq r$ et $k'$ un entier $\leq k$. Si $(U,\varphi,p)$ et $(V,\handwritten{\psi},q)$ sont des cartes de $X$ et $Y$, on a des isomorphismes
\scanpage{44}
% archive sheet number: 43
Revenons aux conditions de la \uline{proposition}, et \unsure{déterminons|déterminer}
\[
\omega_{G_0}=t_{G_0}^{\vee}\simeq M_{\unsure{\nu|v}}/\pi^{\unsure{F|f|r}}M_\nu\simeq \Omega/p\Omega
\]
avec $\deleted{\unsure{K}}k=A/\pi A$ y opérant de la façon évidente, et la structure de vectoriel sur $\inserted{W/pW=}k=\mathbb F_q$ est déduite de la structure précédente grâce à $\sigma^\nu$. Le cas le plus simple est celui où $\nu=0$, i.e. où \deleted{\unsure{les deux}} les deux opérations de $k$ sur $\omega_{G_0}$ coïncident. \deleted{\unsure{Si on sup|Soit}} (les autres cas s'en déduisent grâce à \unsure{autom} $\sigma^\nu$ de $k=\mathbb F_q$ \deleted{\unsure{sur $k_p$}}).

La filtration canonique sur $M\otimes_W k=\sum M_\nu\otimes_A k$ donne comme sous-module (de rang $1$) \deleted{\unsure{$M_{f-1}\otimes_A k$}} \inserted{$\omega_0$ \deleted{\unsure{$=$}} tel que $\omega_0^{(p)}\subset(M\otimes_W k)^{(p)}\unsure{=}$ soit $M_{f-1}\otimes_A k$ ce qui donne $\omega_0=M_0\otimes_A\deleted{\unsure{k}}k$.}

Le cas \unsure{standard|standart} est
\[
G_0=G_0(A,\pi)\quad\text{correspondant à }\;(\nu=0\text{ et })\;\Omega=A
\]
on en déduit le cas d'un $\Omega$ quelconque par
\[
\deleted{\unsure{G_0}}\quad\Omega\otimes_A G_0(A)
\]

Donc on a dans ce cas comme module de Dieudonné
\[
\left\{\begin{aligned}
M&=A\otimes_{\mathbb Z_p}W\\
F_M&=\mathrm{id}_A\otimes_{\mathbb Z_p}\sigma+\operatorname{pr}_0\bigl((\pi-1)\otimes\sigma\bigr)
\end{aligned}\right.
\]
où $\operatorname{pr}_0$ est la projection sur le composant correspondant au $A$-hom $A\otimes_{\mathbb Z_p}W\longrightarrow A$ qui sur $W$ est l'inclusion $W\mathrel{\overset{i}{\unsure{\lhook\joinrel\relbar|\hookrightarrow}}}A$.

\begin{center}\rule{2.2em}{0.3pt}\\[-4pt]\rule{1.2em}{0.3pt}\end{center}

Nous nous proposons maintenant de relever $G_0$ \deleted{l'action} en un groupe de BT $G$ sur $A$, \inserted{comm.} avec action de $A$ sur $G$, de telle façon que l'action de $A$ sur $t_G$ (qui est un $A$-module inversible) soit l'action évidente. D'après Lubin-Tate,
\scanpage{45}
% archive sheet number: 44
\mspage{78}
\hfill n\textsuperscript{o} 382

\[
P^k_x:C^r_x(X)\longrightarrow P^k_x(X)
\]
soit un homomorphisme ; on a les formules :
\[
P^k_x(f)+P^k_x(g)=P^k_x(f+g)\;,\quad\lambda.P^k_x(f)=P^k_x(\lambda.f)\;,\quad P^k_x(f).P^k_x(g)=P^k_x(f.g)\;,
\]
pour $f,g\in C^r_x(X)$, $\lambda\in K$.

7.1.3. Soit de plus $h$ un morphisme dans une variété $Z$ de classe $C^r$ d'un voisinage de $y$. Alors $j^k_x(h\circ f)$ ne dépend que de $j^k_x(f)$ et de $j^k_y(h)$ et on l'appelle le \uline{produit} (ou le \uline{composé}) de $j^k_x(f)$ et $j^k_y(h)$ ; le composé $v\circ u$ de deux jets $u$ et $v$ est donc défini si le but de $u$ coïncide avec la source de $v$ ; ce produit est associatif. Un jet $u$ d'ordre $k$ de $X$ dans $Y$ est \uline{inversible} s'il existe un jet $v\in J^k(Y,X)$ tel que $v\circ u$ et $u\circ v$ soient définis et égaux aux jets des applications identiques $1_X$ et $1_Y$. Si $k=0$, tout jet est inversible ; si $k\geq 1$, un jet $u$ est inversible si et seulement si c'est la partie principale d'ordre $k$ d'un isomorphisme local de $X$ dans $Y$ ; dans ce cas, tout morphisme $f$ tel que $j^k_x(f)=u$ est un isomorphisme local en $x$. Un jet d'ordre $k$ inversible de $X$ dans $K^p$ (resp. de source $x$) est appelé un \uline{repère d'ordre} $k$ de $X$ (resp. de source $x$). L'ensemble de ces repères est noté $R^k(X)$ (resp. $R^k_x(X)$). Si $U$ et $V$ sont des ouverts de $X$ et $Y$, on a des inclusions canoniques
\[
J^k(U,V)\subset J^k(X,Y)\qquad\text{et}\qquad R^k(U)\subset R^k(X)\;.
\]

7.1.4. Si $k'\leq k$, le jet $j^{k'}_x(f)$ d'ordre $k'$ d'un germe de morphisme $f$ ne dépend que de $j^k_x(f)$ et est inversible si $j^k_x(f)$ l'est. En faisant correspondre $j^{k'}_x(f)$ à $j^k_x(f)$, on définit des applications canoniques surjectives
\[
J^k_x(X,Y)_y\longrightarrow J^{k'}_x(X,Y)_y\;,\qquad R^k_x(X)\longrightarrow R^{k'}_x(X)
\]
et $P^k_x(X)\longrightarrow P^{k'}_x(X)$ ; la dernière est un homomorphisme d'algèbres.

7.1.5. Soit $f=(f^1,\ldots,f^q)$ une application de classe $C^r$ d'un ouvert $U$ de $K^p$ dans l'espace $K^q$. Si $x\in U$, le jet $j^k_x(f)$ de $f$ est complètement déterminé par les coordonnées de sa source $x$ et par les valeurs en $x$ des fonctions $\Delta^\alpha f^i$, avec
\scanpage{46}
% archive sheet number: 45
un tel \unsure{relèvt|relèvement} existe et est unique (à isom \inserted{\unsure{(non unique|canonique|unique}} près). \deleted{\illegible{1}} Pour simplifier et pouvoir appliquer la théorie des déformations, on va supposer que $\pi A$ est stable par puissances divisées. Donc remonter $G_0$ en un gpe de B.T $G$ sur $A$ (pour le moment sans nous préoccuper de l'action de $A$) revient à remonter la filtration de $M\otimes_W k=\deleted{\unsure{M}}M_A\otimes_A k$ en une filtration de $M_A$ (sur $A$).

% Left of the following paragraph. The first row consists of three unconnected labels.
\diagram{\begin{tikzcd}[ampersand replacement=\&,column sep=small,row sep=small]
\Omega\&M=\Omega\otimes_{\mathbb Z_p}W\&M_A=\Omega\otimes_{\mathbb Z_p}A\\
A\arrow[r,dash]\arrow[d,dash]\&A\otimes_{\mathbb Z_p}W\arrow[r,dash]\arrow[d,dash]\&A\otimes_{\mathbb Z_p}A\arrow[d,dash]\\
\mathbb Z_p\arrow[r,dash]\&W\arrow[r,dash]\&A
\end{tikzcd}}

La compatibilité à l'action de $A$ signifie que c'est une filtration de $A\otimes_{\mathbb Z_p}A$-modules, i.e. on a un sous-\inserted{$A\otimes_{\mathbb Z_p}A$-}module \deleted{de rang $1$} de $\Omega\otimes_{\mathbb Z_p}A$, \inserted{de rang $1$ sur $1\otimes A$} ce dernier sera \deleted{\unsure{donc}} d'ailleurs can. isom à $\omega_G$

Or
\diagram{\begin{tikzcd}[ampersand replacement=\&,column sep=small]
\unsure{\cdot|0}\arrow[r,dash]\&\omega_G\arrow[r]\&\Omega\otimes_{\mathbb Z_p}A\arrow[r]\&t_{G^*}\arrow[r]\&0
\end{tikzcd}}
\unsure{ou|i.e.|et}, par transposition
\diagram{\begin{tikzcd}[ampersand replacement=\&,column sep=small,row sep=small]
\unsure{(***)|(**)}\&\cdot\arrow[r,dash]\&\omega_{G^*}\arrow[r,dash]\&\Omega^\vee\otimes_{\mathbb Z_p}A\arrow[r,"{\deleted{\unsure{\beta}}\rho}"]\&t_G\arrow[r,dash]\&0\\
\&\&\&A\otimes_{\mathbb Z_p}A\arrow[r,dash,"\rho_0"]\&A\&
\end{tikzcd}}

L'hypothèse que l'action de $A$ \inserted{sur $\omega_G$} (via action sur $G$) est l'action de $A$-module naturelle, signifie que \deleted{\unsure{$\rho$ est compatible|soit compatible}} \inserted{la structure} de $A\otimes_{\mathbb Z_p}A$-module sur $\omega_G$ (ou $t_G$) \deleted{\unsure{avec $\rho$ i.e. $\rho_0$ sur|avec $\rho$ i.e. $p_0$ sur}} provient de sa structure de $A$-module \deleted{triviale} évidente par restriction des scalaires à l'aide de
\[
\rho_0:A\otimes_{\mathbb Z_p}A\mathrel{\unsure{\relbar|\longrightarrow}} A\qquad \rho_0(\lambda)\otimes\mu)=\lambda\mu
\]
\scanpage{47}
% archive sheet number: 46
\mspage{77}
\hfill n\textsuperscript{o} 382

\S\ 7. \uuline{Jets. Opérateurs différentiels}.

A partir de maintenant, on adopte la convention suivante : si $k$ est un entier, et si $r$ désigne l'un des deux symboles $\infty$ et $\omega$, on pose $r-k=r$. Ainsi les dérivées d'ordre $k$ d'une fonction de classe $C^r$ dans un ouvert de $K^n$ sont de classe $C^{r-k}$ dans cet ouvert, que $r$ soit fini ou non.

7.1 - \uline{Jets d'applications}.

7.1.1. Soient $X$ et $Y$ deux variétés de classe $C^r$, et soient $x\in X$, $y\in Y$ ; soient $f$ et $g$ deux morphismes d'un voisinage de $x$ à valeurs dans $Y$, tels que $f(x)=g(x)=y$ ; soit $k$ un entier positif $\leq r$, ou éventuellement le symbole $\infty$ si $r\geq\infty$. On dit que $f$ et $g$ ont un \uline{contact d'ordre} $\geq k$ \uline{en} $x$ s'il existe une carte $(U,\varphi,p)$ de $X$ en $x$ et une carte $(V,\handwritten{\psi},q)$ de $Y$ en $y$ telles que $\handwritten{\psi}\circ f\circ\varphi^{-1}$ et $\handwritten{\psi}\circ g\circ\varphi^{-1}$ aient un contact d'ordre $\geq k$ au point $\varphi(x)\in K^p$ (cf. n\textsuperscript{o} 1.1.1, 2.5.4). Cette propriété est alors vérifiée quel que soit le choix des cartes $(U,\varphi,p)$ et $(V,\handwritten{\psi},q)$ et définit une relation d'équivalence entre germes de morphismes. Une classe d'équivalence pour cette relation est appelée un \uline{jet d'ordre} $k$ \uline{de morphisme de} $X$ \uline{dans} $Y$, \uline{de source} $x$ \uline{et de but} $y$ ; la classe du morphisme $f$ est notée $j^k_x(f)$ et se nomme le \uline{jet d'ordre} $k$ de $f$ en $x$, ou \uline{la partie principale d'ordre} $k$ \uline{de} $f$ en $x$. L'ensemble des jets d'ordre $k$ de $X$ dans $Y$ (resp. de source $x$, \handwritten{\unsure{par|pr}} resp. de but $y$) est noté $J^k(X,Y)$ (resp. $J^k_x(X,Y)$, resp. $J^k(X,Y)_y$), et l'on pose :
\[
J^k_x(X,Y)_y=J^k_x(X,Y)\handwritten{\cap}J^k(X,Y)_y\;.
\]

7.1.2. Si $Y=K$, $J^k(X,Y)$ est l'ensemble des parties principales d'ordre $k$ de fonctions morphiques sur $X$. On écrira en général $P^k(X)$, $P^k_x(X)$ et $P^k_x(f)$ pour $J^k(X,K)$, $J^k_x(X,K)$ et $j^k_x(f)$.

Soit $C^r_x(X)$ l'ensemble des germes en $x$ de fonctions morphiques sur un voisinage de $x$ ; c'est une $K$-algèbre. Il existe sur $P^k_x(X)$ une structure de $K$-algèbre et une seule telle que l'application
\scanpage{48}
% archive sheet number: 47
Cela, plus le fait que $t_G$ et $\Omega^\vee\otimes_{\mathbb Z_p}A$ sont de rang $1$ sur $A$ resp $A\otimes_{\unsure{\mathbb Z_p|\partial_p}}A$, et que $\rho$ est surjectif, prouve que $t_G$ \uline{est déduit de} $\Omega^\vee\otimes_{\mathbb Z_p}A$ par le changt de base $A\otimes_{\mathbb Z_p}A\mathrel{\overset{\rho_0}{\relbar}}A$, donc \deleted{\unsure{par suite}} (comme $\Omega^\vee\otimes_{\mathbb Z_p}A$ est déduit de $\Omega^\vee$ par $A\xrightarrow{\lambda\to\lambda\otimes 1}A\otimes_{\mathbb Z_p}A$) on obtient un isom can.
\[
t_G\simeq\Omega^\vee\qquad\text{i.e.}\qquad\omega_G\simeq\Omega
\]
et la suite exacte $(**)$ est la suite exacte canonique évidente, \inserted{d'où $(*)$ par dualité.} Comme cette filtration étend celle de $M_A\otimes_A k$, on a retrouvé ici le th. d'existence et d'unicité de Lubin-Tate, (\deleted{\unsure{en}}sous une forme moins forte d'ailleurs en ce qui concerne l'unicité, et sous l'hyp. complémentaire que $\exists$ puiss. div. sur $\pi A$.

Si on abandonne cette dernière hypothèse, et qu'on utilise la classification des déformations à isogénie près, on trouve de $=$ qu'on a une suite exacte canonique
\diagram{\begin{tikzcd}[ampersand replacement=\&,column sep=small,row sep=small]
0\arrow[r,dash]\&\omega_G\otimes_A K\arrow[r,dash]\&M_K\arrow[r]\arrow[d,equal]\&t_{G^*}\otimes_{\unsure{A|p}}K\arrow[r,dash]\&0\\
\&\&\Omega\otimes_{\mathbb Z_p}\deleted{\unsure{k}}K\&\&
\end{tikzcd}}
ou encore
\diagram{\begin{tikzcd}[ampersand replacement=\&,column sep=small,row sep=small]
0\arrow[r,dash]\&\omega_{G^*}\otimes_A K\arrow[r,dash]\&\overset{\vee}{M}_K\arrow[r]\arrow[d,equal]\&t_G\otimes_A K\arrow[r]\&0\\
\&\&\Omega^\vee\otimes_{\mathbb Z_p}K\&\&
\end{tikzcd}}
et on trouve encore que
\[
t_G\otimes_A K\simeq\overset{\vee}{M}_K\otimes_{\Omega^\vee\otimes_{\mathbb Z_p}K}(K\otimes K,\rho_0)\simeq\Omega^\vee\otimes_A K
\]
\scanpage{49}
% archive sheet number: 48
\mspage{76}
\hfill n\textsuperscript{o} 382

\[
\begin{aligned}
(f\omega)\wedge\omega'&=f.(\omega\wedge\omega')=\omega\wedge f\omega'\;,\\
i(fs)\omega&=f.i(s)\omega=i(s)(f\omega)\;,\quad\text{etc.}
\end{aligned}
\]

6.8.10. On définit de même la \uline{puissance symétrique} $k$-\uline{ième} $S^k(E)$ d'un espace fibré vectoriel $E$ ; ses fibres sont égales aux puissances symétriques $k$-ièmes des fibres de $E$. On laisse au lecteur le soin de formuler, pour les puissances symétriques des résultats analogues à ceux énoncés ci-dessus pour les puissances extérieures.

6.9 - \uline{Changement de structure}.

6.9.1. Soit $E$ un espace fibré vectoriel de classe $C^r$ ; c'est aussi un espace fibré vectoriel de classe $C^{r'}$ pour tout $r'\leq r$, ainsi qu'un espace fibré vectoriel au sens topologique ; de même, s'il est analytique sur une extension de $K$, il l'est aussi sur $K$.

6.9.2. Les changements de structure définis au n\textsuperscript{o} 6.9.1 sont compatibles avec la formation des images réciproques, ainsi qu'avec le \hspace{3em}procédé de construction du n\textsuperscript{o} 6.6.

\begin{center}-:-:-:-\end{center}
